In this section, the same fitting procedure that are applied to the
  data are applied to a \powheg\ sample of 
  $b\bar{b}$ and $c\bar{c}$ to dimuon decays.
The sample is an official ATLAS MC event-generation sample described 
  in Ref.~\cite{ATLHI-420}.
The same fiducial selections are applied on the single-muons and 
  on the muon-pair as those in the data analysis.
  namely that the muons $|eta|<2.4$ and \pt>4~\GeV, and that 
  for the pair $\ptbar>5$~\GeV\ and $|\Deta|>0.8$.
For the comparisons here, the $b\bar{b}$ and $c\bar{c}$ 
  samples are combined appropriately weighted by the cross-sections.
The samples are generated using nuclear PDFs (nCTEQ15+SIH)~\cite{ATLHI-420}
  corresponding to the Pb-nucleus, and thus include the correct admixture
  of $pp$, $pn$ and $nn$ collisions.
The goal here is not to validate the MC samples, but to get a general idea
  of how close the \powheg\ calculations are to the measurements.

Figure~\ref{fig:powheg_fit} shows the Fits to Eq.~\eqref{eq:fit_func}
  for the \powheg\ sample.
These fits are made constraining the $v_{2,2}$ parameter 0 zero,
  as no meahanism to generate a global elliptic modulation is present
  in this MC sample.
Figure~\ref{fig:powheg_compare} compares the $\Gamma$ and widths 
  obtained in the data analysis to those obtained from the
  \powheg\ samples.
The \powheg\ calculations match with the \pp\ data within a few percent.



\begin{figure}[h]
\centering
\includegraphics[width=1.0\textwidth]{figures/AllFigs/can_dphi_fits_mc3_LorentzZyamNoVn}
\caption{
  Fits to the dimuon correlations using Eq.~\eqref{eq:fit_func}.
  The left, center and right panels correspond to 
    same-sign, opposite-sign and all-pairs respectively.
\label{fig:powheg_fit}
} 
\end{figure}

\begin{figure}[h]
\centering
\includegraphics[width=1.0\textwidth]%
{figures/AllFigs/can_systPbPbPPMC_Tau1_Tight_pTBar_GT5_EffCor}
\includegraphics[width=1.0\textwidth]%
{figures/AllFigs/can_systPbPbPPMC_Width1_Tight_pTBar_GT5_EffCor}
\caption{
  Comparison of the $\Gamma$ (top) and widths obtained from
    the \powheg\ samples to those obtained in the data.
  The left, center and right panels correspond to 
    same-sign, opposite-sign and all-pairs respectively.
\label{fig:powheg_compare}
} 
\end{figure}


Additionally, Figure~\ref{fig:powheg_bb_cc_combined} compares the 
  relative contributions of the \bbbar\ and \ccbar\ samples
  to the combined sample to get an idea of the origin of the
  dimuon pairs.
The correlations measured in the data for \pp\ collisions are also shown.
\begin{figure}[h]
  \centering
  \includegraphics[width=1.0\textwidth]%
  {figures/AllFigs/can_dphi_powheg_mc}
  \caption{
    Comparison of the relative contribution of \bbbar\ and \ccbar\ pairs.
    The left, center and right panels correspond to 
      same-sign, opposite-sign and all-pairs respectively.
    The correlations measured in the data for \pp\ collisions are also shown.
  \label{fig:powheg_bb_cc_combined}
  } 
\end{figure}

\clearpage
\section{\Dphi\ smearing from decays}
\label{sec:powheg_decays}
In this analysis we compare the angular smearing of muon pairs
  from decays of HF hadron pairs.
It is useful to get an idea oh how the $\Dphi$ distributions
  of the muons compare to that of the parent hadrons.
Note that the calculations presented here are to be treated
  as ``back of the envelope calculations'', not as 
  rigorous calculations.
Also, since the results in this analysis compare the \pp\ 
  and \PbPb\ measurements, any smearing from the decay
  produces identical effects in the \pp\ and \PbPb\ measurements
  and drops out when the quadrature difference in the measured widths
  are evaluated.
  

The \powheg\ sample described in Section~\ref{sec:powheg}
  is used to obtain the angular correlations between
  the muon and the last $b$-hadron (or $c$-hadron) in the decay chain.
The distributions are shown in Figure~\ref{fig:powheg_decays}.
The $b\rightarrow\mu$ decay introduces a smearing of $\Dphi\sim0.16$
  while the $c\rightarrow\mu$ decay introduces a smearing of $\Dphi\sim0.06$.
The final measured widths for the combined-sign pairs are $\approx0.67$ in the \PbPb\ data.
This would mean that the width for the original corelation between the hadrons (when both hadrons are $b$)
  would be:
\begin{eqnarray}
  \text{Width}^{h-h}=\sqrt{0.67^2-2*(0.16^2)}\approx0.63,
\end{eqnarray}
where a factor of two is included in the subtraction to account
  for each of the $b$-hadron decays.


\begin{figure}[h]
  \centering
  \includegraphics[width=1.0\textwidth]%
  {figures/AllFigs/can_decay_widths}
  \caption{
    Left panel: distribution of the $\Dphi$ angle between the muon and the 
      last $b$-hadron in the decay chain obtained from the \powheg\ sample.
    Right panel: similar plot for muons from charm-hadron decays.
    The r.m.s. widths of the distributions are listed on the plots.
  \label{fig:powheg_decays}
  } 
\end{figure}
